\documentclass[11pt]{article}
\usepackage[utf8]{inputenc}
\usepackage[T1]{fontenc}
\usepackage{babel}
\usepackage{lmodern}
\usepackage{amsmath,amssymb}
\usepackage{hyperref}
\usepackage{mathptmx}

% THEOREM CONFIGS (ENGLISH)
\usepackage[thmmarks]{ntheorem}

\theoremstyle{plain}
\theoremheaderfont{\normalfont\bfseries}
\theorembodyfont{\normalfont}
\theoremseparator{:}
\theorempreskip{\topsep}
\theorempostskip{\topsep}

\theoremsymbol{\ensuremath{\bigcirc}}
\theorembodyfont{\normalfont}
\newtheorem{definition}{Definition}[section]

\theoremsymbol{\ensuremath{\triangle}}
\theorembodyfont{\normalfont}
\newtheorem{note}[definition]{Note}

\theoremsymbol{}
\theorembodyfont{\normalfont}
\newtheorem{proposition}[definition]{Proposition}

\theoremsymbol{\ensuremath{\square}}
\theorembodyfont{\normalfont}
\newtheorem{freeprop}[definition]{Proposition}

\theoremstyle{nonumberplain}
\theoremheaderfont{\normalfont\bfseries}
\theorembodyfont{\normalfont}
\theoremsymbol{\ensuremath{\blacksquare}}
\newtheorem{proof}[definition]{Proof}

% Generic Macros (english)
\newcommand{\N}{\mathbb{N}}
\newcommand{\Z}{\mathbb{Z}}
\newcommand{\Q}{\mathbb{Q}}
\newcommand{\R}{\mathbb{R}}
\newcommand{\C}{\mathbb{C}}
\newcommand{\set}[1]{\{#1\}}
\newcommand{\inv}[1]{#1^{-1}}

\title{Endomorphisms}
\author{Gian Carlo}
\date{\today}

\begin{document}
\maketitle

% Specific macros

\section{Introduction}

This should (eventually) read like a taxonomy of the elements of the endomorphism of some set. If cardinality arguments appear and look weird, assume the underlying set is finite.

\section{Investigation}

So the first thing I want to lock is notation. For a given set $S$, the set of functions $S \to S$ will be written as $E_S$. It is not pretty but it is easy to type. If we ever need to think about what a given $f \in E_S$ does to some $s \in S$, we shall not write it as $f(s)$ but instead we will compose the constant function $f_s(x) = s$ with $f$. Immediately we will abuse notation and say that $s$ denotes both the element of $S$ and the constant function.

Another thing we will do is write function products as follows \begin{equation}
    fg = g \circ f.
\end{equation}
This of course generalizes to longer products like $fghgfhfhg$, which are associative since composition is associative. Then, instead of evaluating functions (if we need to work with element traces/paths) we will merely write \begin{equation}
    s_0f
\end{equation}
where $s_0$ is the constant function $s_0(x) = s_0 \in S$ and $f$ is some element of $E_S$. We will name set elements after the set, so if $S$ is a set the symbols $s_0, s_1, ..., s_n, ...$ are both constant functions and elements of $S$. For a set $C$, we reserve $c_0, ...$ and so on.

Finally, when speaking about cardinalities, it probably safer to assume finite sets, so we have the common arithmetic available. I know wonder what would take for set theory to allow for arbitrary cardinalities. What is a set of size $2.5$?

\section{Synthesis}

\begin{definition}
    Let $S$ be a set. The set of functions $f : S \to S$ will be called the \textbf{endomorphism} set, or the \textbf{function set}, and will be written as $E_S$.
\end{definition}

\begin{definition}[Function Product] \label{def:function_product}
    Let $S$ be a set and $E_S$ it's function set. The product of two functions $f, g \in E_S$ is defined as their composition: \begin{equation}
        fg = g \circ f.
    \end{equation}
\end{definition}

\begin{note}
    The previous definition is done in such a way as that reading left-to-right also tells us the order of application of the functions.
\end{note}

\begin{proposition}
    The set of endomorphisms of some set $S$ with the product of \ref{def:function_product} forms a monoid.
\end{proposition}

\begin{proof}
    Because the product is defined from the composition, it inherits the associativity and the neutral element $i(x) = x$, the identity function.

    Nevertheless, we show some calculations \footnote{These are calculations I actually understand, so I might as well show them.}. Let $f, g, h \in E_S$ and consider $f(gh)$. By definition, \begin{equation}
        f(gh) = (gh) \circ f
    \end{equation}
    which yields \begin{equation}
        (gh) \circ f = (h \circ g) \circ f
    \end{equation}
    which we associate freely, since it's in composition form, to \begin{equation}
        (h \circ g) \circ f = h \circ (g \circ f),
    \end{equation}
    giving us \begin{equation}
        h \circ (g \circ f) = h \circ (fg) = (fg)h
    \end{equation}
    as expected.

    Consider the function $i(x) = x$ in $E_S$. By: definition, composition property, composition property, and definition, we have $fi = i \circ f = f = f \circ i = if$.
\end{proof}

\begin{note}
    We will not justify every equality like in the last proposition. Do not worry.

    The next proposition serves to establish our notation.
\end{note}

\begin{proposition}[No More Parentheses]
    Let $S$ be a set, $E_S$ it's function set and $\bar s(x) = s \in S$ be a constant function. For every function $f \in E_S$ we have \begin{equation}
        \bar s f = f(s).
    \end{equation}
\end{proposition}

\begin{proof}
    We just apply our definition. Since $\bar s(x) = s$, when we compute $\bar s f$ we obtain \begin{equation}
        \bar s f = f \circ \bar s
    \end{equation}
    which applied to some $x \in S$ yields \begin{equation}
        (f \circ \bar s) (x) = f(\bar s (x)) = f(s)
    \end{equation}
    as expected.
\end{proof}

\begin{note}
    The previous argument exists for the sole purpose of freeing us from having to write $f(s)$ for $s \in S$ should we ever need to think about where an element goes. We will also abuse the notation and consider $s \in S$ both as an element and as it's constant function $s(x) = s$.
\end{note}

\begin{proposition}[Rightmost Zeroes]
    For a given set $S$ and it's endomorphisms, any constant function $s \in E_S$ is such that $fs = s$ for all $f \in E_S$.
\end{proposition}

\begin{proof}
    If we expand the product into the composition \begin{equation}
        fs = s \circ f
    \end{equation}
    and remember $s$ is a constant function, then \begin{equation}
        s \circ f = s 
    \end{equation}
    and then $fs = s$ as expected.
\end{proof}

\begin{note}
    Two things we ought to point out now: \begin{enumerate}
        \item We will write the identity function $f(x) = x$ as $e \in E_S$.
        \item The last two propositions (and others to come) remain true if instead of a single function we compose with a finite product of functions \begin{equation}
            P = \prod_{i = 1}^{n \in \N} f_i.
        \end{equation}
        In other words, $sP$ is still equal to evaluating $P(s)$ and $Ps = s$.
    \end{enumerate}
\end{note}

\section{Not Conclusions}

Next step is actually doing some taxonomy on $E_S$. We have bijections, constant functions. I know the literature calls the elements of the function set transformations of the set but I will instead call them $k$-classifications based on the size of their image, because I will actually use the $k$ in an equation.

I am writing this alone. Why do I use plural language persons? We don't know.

\end{document}